\documentclass[conference]{IEEEtran}

% ---- Packages ----
\usepackage{cite}
\usepackage{amsmath,amssymb,amsfonts}
\usepackage{algorithmic}
\usepackage{graphicx}
\usepackage{textcomp}
\usepackage{xcolor}
\usepackage{booktabs}
\usepackage{multirow}
\usepackage{hyperref}
\usepackage{subcaption}
\usepackage{array}
\usepackage{tabularx}
\usepackage{soul} % for highlights
\usepackage{balance}

% ---- Mock result helper ----
\newcommand{\mock}[1]{{\color{blue}#1}}
\newcommand{\TODO}[1]{{\color{red}\textbf{[TODO: #1]}}}

\begin{document}

\title{Contrastive Curriculum Learning for Sentiment Reasoning in Healthcare: Addressing Keyword Bias and Class Ambiguity}

\author{
\IEEEauthorblockN{Author One\IEEEauthorrefmark{1},
Author Two\IEEEauthorrefmark{2},
Author Three\IEEEauthorrefmark{3}}
\IEEEauthorblockA{\IEEEauthorrefmark{1}Affiliation One, Country\\
\IEEEauthorrefmark{2}Affiliation Two, Country\\
\IEEEauthorrefmark{3}Affiliation Three, Country\\
\{author1, author2, author3\}@email.edu}
}

\maketitle

% ============================================================
% ABSTRACT
% ============================================================
\begin{abstract}
Sentiment analysis in healthcare requires not only accurate classification but also transparent reasoning to support clinical decision-making. Recent work on Sentiment Reasoning introduced a multimodal framework where models predict both a sentiment label and a free-text rationale from medical transcripts. However, current approaches suffer from keyword shortcutting---where surface-level terms override contextual understanding---and systematic confusion between NEUTRAL and polar sentiment classes. In this work, we propose \textbf{Contrastive Curriculum Sentiment Reasoning (CCSR)}, a novel framework that addresses these limitations through three mechanisms: (1) a sentiment-contrastive representation module that enforces class separation in embedding space with rationale-grounded sub-clustering, (2) a difficulty-aware curriculum that progressively introduces ambiguous boundary samples, and (3) focal loss to mitigate class imbalance. On the English subset of the Sentiment Reasoning benchmark, CCSR achieves \mock{70.23\%} accuracy and \mock{69.15\%} macro-F1, representing a \mock{+5.69\%} and \mock{+4.79\%} absolute improvement over the previous best. Rationale quality also improves, with BERTScore reaching \mock{0.847} and a new faithfulness metric showing \mock{87.3\%} of generated rationales logically entail their predicted labels. Ablation studies confirm that the contrastive objective contributes the majority of gains, validating our thesis that representation quality---not model scale alone---is the key bottleneck in healthcare sentiment reasoning.
\end{abstract}

\begin{IEEEkeywords}
Sentiment analysis, healthcare NLP, contrastive learning, curriculum learning, rationale generation, explainable AI
\end{IEEEkeywords}

% ============================================================
% 1. INTRODUCTION
% ============================================================
\section{Introduction}

Sentiment analysis plays a pivotal role in healthcare, enabling real-time evaluation of patient satisfaction~\cite{xia2009improving}, monitoring emotional well-being~\cite{cambria2012sentic}, and supporting mental health assessment~\cite{pestian2012sentiment}. As healthcare increasingly integrates AI-driven decision support, the transparency of these systems becomes critical---errors can have severe consequences, and clinicians need to understand \emph{why} a model made a particular prediction~\cite{antoniadi2021current}.

Nguyen et al.~\cite{nguyen2024sentiment} addressed this need by introducing \emph{Sentiment Reasoning}, a novel task where models predict both a sentiment label (POSITIVE, NEUTRAL, or NEGATIVE) and generate a free-text rationale explaining the prediction. Their work established the world's largest multimodal sentiment analysis dataset (30,000 samples across five languages) built on real-world doctor-patient conversations from VietMed~\cite{leduc2024vietmed}. They demonstrated that rationale-augmented training improves classification by approximately 2\% in both accuracy and macro-F1.

However, their error analysis reveals fundamental limitations in current approaches. As shown in their confusion matrix analysis, models systematically misclassify POSITIVE and NEGATIVE transcripts as NEUTRAL at rates of 23--27\%. Our investigation identifies three root causes:

\begin{enumerate}
    \item \textbf{Keyword shortcutting}: Models latch onto sentiment-bearing keywords (e.g., ``helpful,'' ``diabetes'') rather than contextual reasoning. Their error analysis (Table 9 in~\cite{nguyen2024sentiment}) shows high-confidence misclassifications driven by isolated keywords overriding the actual sentiment.
    \item \textbf{No explicit class separation}: The generation loss alone provides no signal for where class boundaries lie in representation space, leaving the model to discover them implicitly.
    \item \textbf{Uniform training}: Ambiguous boundary samples---where sentiment is subtle or context-dependent---receive the same training treatment as clear-cut examples, diluting the learning signal on the hardest cases.
\end{enumerate}

To address these challenges, we propose \textbf{Contrastive Curriculum Sentiment Reasoning (CCSR)}, a framework that jointly optimizes sentiment classification, rationale generation, and representation quality. Our contributions are:

\begin{enumerate}
    \item A novel \textbf{CCSR framework} that combines a sentiment-contrastive representation module with curriculum-based training for joint classification and rationale generation.
    \item A \textbf{difficulty-aware curriculum strategy} that scores sample difficulty using keyword-label mismatch and classifier uncertainty, progressively introducing harder examples.
    \item A \textbf{rationale-grounded contrastive objective} that creates sub-cluster structure within sentiment classes based on reasoning patterns, improving rationale diversity and faithfulness.
    \item \textbf{State-of-the-art results} on the English subset of the Sentiment Reasoning benchmark, with significant improvements in both classification metrics and rationale quality, along with two new evaluation metrics: rationale faithfulness and keyword reliance reduction.
\end{enumerate}

% ============================================================
% 2. RELATED WORK
% ============================================================
\section{Related Work}

\subsection{Sentiment Analysis in Healthcare}

Sentiment analysis in healthcare spans medical web mining~\cite{ali2013can, na2012sentiment, sokolova2013joe}, biomedical literature analysis~\cite{niu2005analysis, sarker2011outcome}, and clinical text processing~\cite{pestian2012sentiment, cambria2012sentic}. While these works address text-based sentiment, speech-based sentiment analysis in healthcare remains underexplored. Nguyen et al.~\cite{nguyen2024sentiment} were the first to introduce speech sentiment analysis specifically within the healthcare domain, establishing the Sentiment Reasoning task and dataset.

\subsection{Rationale-Augmented Training and CoT Distillation}

Chain-of-Thought (CoT) distillation has shown that training generative models on rationale-augmented targets improves performance. Hsieh et al.~\cite{hsieh2023distilling} proposed Distilling Step-by-Step for encoder-decoders, while Chen et al.~\cite{chen2024post} demonstrated that post-thinking (label before rationale) is more stable and token-efficient than pre-thinking. Wadhwa et al.~\cite{wadhwa2024investigating} further investigated CoT-augmented distillation, finding that rationale format has limited impact---a finding confirmed by Nguyen et al. Our work builds on post-thinking but adds contrastive and curriculum signals that go beyond the rationale format question.

\subsection{Contrastive Learning for NLP}

Supervised contrastive learning (SupCon)~\cite{khosla2020supervised} has been successfully applied to NLP tasks. SimCSE~\cite{gao2021simcse} showed that contrastive objectives dramatically improve sentence embeddings. In sentiment analysis, contrastive learning has been used to improve aspect-based classification~\cite{li2021learning} and cross-domain transfer~\cite{wu2022contrastive}. However, no prior work has combined contrastive learning with generative rationale production in a joint framework. Our sentiment-contrastive module is novel in that it operates on decoder hidden states at a specific generation point and incorporates rationale-grounded sub-clustering.

\subsection{Curriculum Learning for NLP}

Curriculum learning~\cite{bengio2009curriculum} trains models on progressively harder examples. In NLP, it has improved machine translation~\cite{platanios2019competence}, relation extraction~\cite{zhang2019curriculum}, and text classification~\cite{xu2020curriculum}. For sentiment analysis specifically, curriculum approaches have addressed noisy labels~\cite{wang2021curriculum} and class imbalance~\cite{zhou2020curriculum}. Our curriculum strategy is distinguished by its difficulty scoring function, which combines classifier uncertainty with keyword-label mismatch---a measure designed specifically for the keyword shortcutting problem in medical sentiment.

% ============================================================
% 3. METHOD
% ============================================================
\section{Method}

\subsection{Problem Formulation}

Following Nguyen et al.~\cite{nguyen2024sentiment}, let $w_1^N = w_1, w_2, \ldots, w_N$ be a transcript of length $N$ and $C = \{\text{NEGATIVE}, \text{NEUTRAL}, \text{POSITIVE}\}$ be the set of sentiment classes. The Sentiment Reasoning task requires building a model $f$ that both predicts a sentiment class:
\begin{equation}
    w_1^N \rightarrow \hat{c} = \arg\max_{c \in C} f(c | w_1^N)
\end{equation}
and generates a rationale sequence $r_1^M$ of length $M$:
\begin{equation}
    w_1^N \rightarrow r_1^M = \arg\max_{r^*} h(r^* | w_1^N)
\end{equation}

In the post-thinking formulation, both tasks are unified into a single generation target: given input transcript $w_1^N$, the model generates the sequence $[\hat{c}; r_1^M]$ (label followed by rationale).

\subsection{Framework Overview}

Figure~\ref{fig:framework} illustrates the CCSR framework. It consists of three components built on a shared Llama~3 8B backbone: (A) a generative core for joint classification and rationale generation, (B) a sentiment-contrastive representation module, and (C) a difficulty-aware curriculum scheduler.

\begin{figure}[t]
    \centering
    \fbox{\parbox{0.95\columnwidth}{
        \centering
        \vspace{2em}
        \textbf{CCSR Framework Overview}\\[1em]
        \TODO{Insert framework diagram showing:\\
        - Input transcript flowing into Llama 3 backbone\\
        - Component A: Generation head producing Label + Rationale\\
        - Component B: Projection head at end-of-label token $\rightarrow$ SupCon loss\\
        - Component C: Curriculum scheduler controlling data flow\\
        - Joint loss: $\mathcal{L}_{total} = \mathcal{L}_{gen} + \lambda_{con}\mathcal{L}_{con} + \lambda_{bal}\mathcal{L}_{focal}$}
        \vspace{2em}
    }}
    \caption{Overview of the Contrastive Curriculum Sentiment Reasoning (CCSR) framework. Given a medical transcript, the shared backbone produces hidden representations that feed three objectives: generative cross-entropy for label and rationale production (Component A), supervised contrastive loss on projected embeddings at the end-of-label position (Component B), and focal loss for class-balanced classification (integrated via Component C's curriculum schedule). At inference, only Component A is active.}
    \label{fig:framework}
\end{figure}

\subsection{Component A: Generative Core}

The generative core follows the post-thinking approach~\cite{chen2024post, nguyen2024sentiment}. Given transcript $w_1^N$, the model generates the target sequence $y = [\hat{c}; r_1^M]$ autoregressively. The training loss is the standard cross-entropy:
\begin{equation}
    \mathcal{L}_{gen} = -\sum_{t=1}^{|y|} \log p_\theta(y_t | y_{<t}, w_1^N)
\end{equation}

The input is formatted as an instruction:
\begin{quote}
    \texttt{Classify the sentiment and provide a rationale: \{transcript\}}
\end{quote}
with the target:
\begin{quote}
    \texttt{\{LABEL\} \{RATIONALE\}}
\end{quote}

\subsection{Component B: Sentiment-Contrastive Module}

The key insight behind this component is that the generation loss $\mathcal{L}_{gen}$ alone provides no explicit signal for class boundary learning in representation space. To address this, we introduce a contrastive objective operating on the model's internal representations.

\subsubsection{Projection Head}

Let $\mathbf{h}_i \in \mathbb{R}^d$ be the hidden state of the backbone at the \emph{end-of-label token position} for sample $i$---the point where the model has committed to a sentiment class but before rationale generation begins. We project this through a 2-layer MLP:
\begin{equation}
    \mathbf{z}_i = g(\mathbf{h}_i) = W_2 \cdot \text{ReLU}(W_1 \cdot \mathbf{h}_i + b_1) + b_2
\end{equation}
where $W_1 \in \mathbb{R}^{256 \times d}$, $W_2 \in \mathbb{R}^{128 \times 256}$, and $\mathbf{z}_i$ is $\ell_2$-normalized.

\subsubsection{Supervised Contrastive Loss}

For a batch of $B$ samples with labels $\{c_i\}_{i=1}^B$, the supervised contrastive loss is:
\begin{equation}
    \mathcal{L}_{con} = \sum_{i=1}^{B} \frac{-1}{|P(i)|} \sum_{p \in P(i)} \log \frac{\exp(\mathbf{z}_i \cdot \mathbf{z}_p / \tau)}{\sum_{j \neq i} \exp(\mathbf{z}_i \cdot \mathbf{z}_j / \tau)}
\end{equation}
where $P(i) = \{j : c_j = c_i, j \neq i\}$ is the set of positive pairs (same class) and $\tau = 0.07$ is the temperature parameter.

\subsubsection{Rationale-Grounded Sub-Clustering}

Within each sentiment class, we further encourage samples with semantically similar rationales to cluster together. We precompute a rationale similarity matrix $S$ where $S_{ij} = \text{BERTScore}(r_i, r_j)$ for samples $i, j$ sharing the same label. For pairs with $S_{ij} > 0.85$, we add a soft attraction term:
\begin{equation}
    \mathcal{L}_{sub} = \sum_{(i,j): S_{ij} > \delta} S_{ij} \cdot \| \mathbf{z}_i - \mathbf{z}_j \|^2
\end{equation}
where $\delta = 0.85$. This is folded into $\mathcal{L}_{con}$ with a weighting factor $\alpha_{sub} = 0.1$.

The sub-clustering creates meaningful internal structure: within NEGATIVE, for example, ``disease progression'' samples cluster separately from ``emotional distress'' samples, improving the model's ability to generate class-appropriate and diverse rationales.

\subsection{Component C: Difficulty-Aware Curriculum}

\subsubsection{Difficulty Scoring}

We assign each training sample a difficulty score $d_i \in [0, 1]$ computed from two signals:

\textbf{Classifier uncertainty}: We train a lightweight BERT-base classifier on the English subset for 3 epochs and record the predictive entropy for each sample:
\begin{equation}
    u_i = -\sum_{c \in C} p(c | w_i) \log p(c | w_i)
\end{equation}

\textbf{Keyword-label mismatch}: We construct a sentiment lexicon $L$ from training rationales using TF-IDF, identifying words that appear disproportionately in POSITIVE or NEGATIVE rationales. For each transcript $w_i$ with label $c_i$, we compute:
\begin{equation}
    m_i = \frac{|\{w \in w_i : w \in L \wedge \text{sent}(w) \neq c_i\}|}{|\{w \in w_i : w \in L\}| + \epsilon}
\end{equation}
where $\text{sent}(w)$ is the dominant sentiment of word $w$ in the lexicon.

The combined difficulty score is:
\begin{equation}
    d_i = 0.6 \cdot \tilde{u}_i + 0.4 \cdot \tilde{m}_i
\end{equation}
where $\tilde{u}_i$ and $\tilde{m}_i$ are min-max normalized values.

\subsubsection{Curriculum Schedule}

Training proceeds in three phases:
\begin{itemize}
    \item \textbf{Phase 1} (epochs 1--2): $d_i < 0.4$ ($\sim$55--60\% of data)
    \item \textbf{Phase 2} (epochs 3--4): $d_i < 0.7$ ($\sim$85\% of data)
    \item \textbf{Phase 3} (epochs 5--7): All samples, with class-balanced sampling using inverse frequency weights
\end{itemize}

\subsection{Joint Training Objective}

The overall loss combines all three components:
\begin{equation}
    \mathcal{L}_{total} = \mathcal{L}_{gen} + \lambda_{con} \cdot \mathcal{L}_{con} + \lambda_{bal} \cdot \mathcal{L}_{focal}
    \label{eq:total_loss}
\end{equation}
where $\mathcal{L}_{focal}$ is the focal loss~\cite{ross2017focal} with $\gamma = 2$ applied to the classification token, $\lambda_{con} = 0.3$, and $\lambda_{bal} = 0.5$.

At inference time, Component B (projection head and contrastive loss) is discarded entirely. The model generates $[\hat{c}; r_1^M]$ via standard autoregressive decoding, with zero architectural overhead compared to the SFT baseline.

% ============================================================
% 4. EXPERIMENTAL SETUP
% ============================================================
\section{Experimental Setup}

\subsection{Dataset}

We evaluate on the \textbf{English subset} of the Sentiment Reasoning dataset~\cite{nguyen2024sentiment}, which contains 5,695 training and 2,183 test samples. The class distribution is imbalanced: NEUTRAL ($\sim$50\%), NEGATIVE ($\sim$30\%), and POSITIVE ($\sim$20\%). Each sample consists of a transcript (translated from Vietnamese medical conversations), a sentiment label, and a human-written rationale.

\subsection{Baselines}

We compare against the models reported by Nguyen et al. and add our backbone ablation:

\begin{itemize}
    \item \textbf{Encoders (label only)}: mBERT~\cite{devlin2018bert}, BERT~\cite{devlin2018bert}
    \item \textbf{Encoder-decoders}: mT0~\cite{muennighoff2022crosslingual}, Flan-T5~\cite{chung2022scaling} (label only and label + rationale)
    \item \textbf{Decoders}: Mistral-7B~\cite{jiang2023mistral} (label only and label + rationale)
    \item \textbf{Our backbone}: Llama~3 8B-Instruct with SFT only (no CCSR components)
\end{itemize}

\subsection{Ablation Variants}

To isolate each component's contribution, we evaluate seven variants (Table~\ref{tab:ablation}).

\subsection{Training Configuration}

We fine-tune Llama~3 8B-Instruct using LoRA~\cite{hu2021lora} ($r = 16$, $\alpha = 32$) on all attention layers. We use AdamW with learning rate $2 \times 10^{-4}$ and cosine schedule with 5\% warmup. Batch size is 32 across GPUs using DeepSpeed ZeRO Stage 2. Training runs for 7 epochs with early stopping (patience = 2) on validation macro-F1. All experiments are run with 3 random seeds; we report mean $\pm$ standard deviation.

\subsection{Evaluation Metrics}

\textbf{Classification}: Accuracy, class-wise F1 (NEG, NEU, POS), and macro-F1.

\textbf{Rationale quality}: ROUGE-1, ROUGE-2, ROUGE-L, ROUGE-Lsum, and BERTScore.

\textbf{New metrics}:
\begin{itemize}
    \item \textbf{Faithfulness}: We use a DeBERTa-v3 model fine-tuned on MultiNLI to check whether each generated rationale entails the predicted label. We report the percentage of rationales classified as ``entailment.''
    \item \textbf{Keyword reliance}: We compute the average attention entropy over input tokens. Higher entropy indicates less concentration on individual keywords, suggesting more contextual reasoning.
\end{itemize}

Statistical significance is assessed via Student's $t$-test ($\alpha = 0.05$) across seeds.

% ============================================================
% 5. RESULTS AND ANALYSIS
% ============================================================
\section{Results and Analysis}

\subsection{Main Results}

Table~\ref{tab:main_results} presents the performance of all models on the English human transcript.

\begin{table*}[t]
\centering
\caption{Performance on the English human transcript. \textbf{Bold} = best result. $\dagger$ = results from Nguyen et al.~\cite{nguyen2024sentiment}. \mock{Blue values} are mock results pending experiments.}
\label{tab:main_results}
\begin{tabular}{llcccccccccc}
\toprule
\textbf{Model} & \textbf{Setting} & \textbf{Acc.} & \textbf{F1\textsubscript{NEG}} & \textbf{F1\textsubscript{NEU}} & \textbf{F1\textsubscript{POS}} & \textbf{Mac-F1} & \textbf{R-1} & \textbf{R-2} & \textbf{R-L} & \textbf{BERTsc.} & \textbf{Faith.} \\
\midrule
\multicolumn{12}{l}{\textit{Encoder (Label Only)}} \\
mBERT$^\dagger$ & Label & 0.6001 & 0.5972 & 0.6320 & 0.5408 & 0.5900 & -- & -- & -- & -- & -- \\
BERT$^\dagger$ & Label & 0.6143 & 0.6338 & 0.6245 & 0.5653 & 0.6079 & -- & -- & -- & -- & -- \\
\midrule
\multicolumn{12}{l}{\textit{Encoder-Decoder (Label Only)}} \\
mT0$^\dagger$ & Label & 0.6216 & 0.6303 & 0.6418 & 0.5670 & 0.6130 & -- & -- & -- & -- & -- \\
Flan-T5$^\dagger$ & Label & 0.6157 & 0.6295 & 0.6385 & 0.5462 & 0.6048 & -- & -- & -- & -- & -- \\
\midrule
\multicolumn{12}{l}{\textit{Encoder-Decoder (Label + Rationale)}} \\
mT0$^\dagger$ & +Rat. & 0.6175 & 0.6495 & 0.6253 & 0.5535 & 0.6094 & 0.3571 & 0.2202 & 0.3350 & 0.8044 & -- \\
Flan-T5$^\dagger$ & +Rat. & 0.6326 & 0.6487 & 0.6390 & 0.5978 & 0.6285 & 0.3956 & 0.2652 & 0.3728 & 0.8106 & -- \\
\midrule
\multicolumn{12}{l}{\textit{Decoder (Label Only)}} \\
Mistral-7B$^\dagger$ & Label & 0.6290 & 0.6536 & 0.6322 & 0.5850 & 0.6236 & -- & -- & -- & -- & -- \\
\midrule
\multicolumn{12}{l}{\textit{Decoder (Label + Rationale)}} \\
Mistral-7B$^\dagger$ & +Rat. & 0.6454 & 0.6768 & 0.6364 & 0.6176 & 0.6436 & 0.3853 & 0.2386 & 0.3663 & 0.8092 & -- \\
Llama~3 8B & +Rat. & \mock{0.6671} & \mock{0.6923} & \mock{0.6578} & \mock{0.6384} & \mock{0.6628} & \mock{0.4091} & \mock{0.2573} & \mock{0.3842} & \mock{0.8215} & \mock{0.781} \\
\midrule
\multicolumn{12}{l}{\textit{Decoder (CCSR -- Ours)}} \\
\textbf{Llama~3 8B + CCSR} & +Rat. & \mock{\textbf{0.7023}} & \mock{\textbf{0.7268}} & \mock{\textbf{0.6917}} & \mock{\textbf{0.6661}} & \mock{\textbf{0.6915}} & \mock{\textbf{0.4532}} & \mock{\textbf{0.2914}} & \mock{\textbf{0.4187}} & \mock{\textbf{0.8473}} & \mock{\textbf{0.873}} \\
\bottomrule
\end{tabular}
\end{table*}

CCSR achieves \mock{70.23\%} accuracy and \mock{69.15\%} macro-F1, representing absolute improvements of \mock{+5.69\%} and \mock{+4.79\%} over the previous best (Mistral-7B with rationale). Compared to our own SFT-only baseline (Llama~3 8B without CCSR), the gains are \mock{+3.52\%} accuracy and \mock{+2.87\%} macro-F1, demonstrating that the improvements stem from our method, not merely from using a stronger backbone.

The POSITIVE class---historically the weakest due to class imbalance---sees the largest relative F1 improvement (\mock{+4.85\%} over Mistral-7B), confirming that the combination of focal loss and curriculum-based class balancing effectively addresses the underrepresentation.

\subsection{Ablation Study}

\begin{table}[t]
\centering
\caption{Ablation study on the English human transcript. Con = Contrastive, Cur = Curriculum, Foc = Focal Loss. \mock{Blue values} are mock results.}
\label{tab:ablation}
\begin{tabular}{lccc|cc}
\toprule
\textbf{Variant} & \textbf{Con} & \textbf{Cur} & \textbf{Foc} & \textbf{Acc.} & \textbf{Mac-F1} \\
\midrule
SFT-only & & & & \mock{0.6671} & \mock{0.6628} \\
+ Focal & & & \checkmark & \mock{0.6724} & \mock{0.6693} \\
+ Contrastive & \checkmark & & & \mock{0.6893} & \mock{0.6827} \\
+ Curriculum & & \checkmark & & \mock{0.6762} & \mock{0.6731} \\
+ Con + Foc & \checkmark & & \checkmark & \mock{0.6942} & \mock{0.6871} \\
+ Cur + Foc & & \checkmark & \checkmark & \mock{0.6813} & \mock{0.6768} \\
\textbf{CCSR (full)} & \checkmark & \checkmark & \checkmark & \mock{\textbf{0.7023}} & \mock{\textbf{0.6915}} \\
\bottomrule
\end{tabular}
\end{table}

Table~\ref{tab:ablation} reveals the marginal contribution of each component. The contrastive module alone contributes the largest gain (\mock{+1.99\%} macro-F1 over SFT-only), confirming our thesis that explicit class separation in representation space is the primary bottleneck. The curriculum adds \mock{+1.03\%}, and focal loss adds \mock{+0.65\%}. Crucially, the full combination outperforms any pair, indicating complementary effects: contrastive learning improves representation quality, curriculum learning sharpens boundary discrimination, and focal loss rebalances the class prior.

\subsection{Confusion Matrix Analysis}

\begin{figure}[t]
    \centering
    \begin{subfigure}[b]{0.48\columnwidth}
        \centering
        \fbox{\parbox{0.9\columnwidth}{
            \vspace{1.5em}
            \centering
            \TODO{Confusion matrix:\\Mistral-7B + Rationale\\(from paper or reproduced)}
            \vspace{1.5em}
        }}
        \caption{Mistral-7B + Rationale}
        \label{fig:cm_baseline}
    \end{subfigure}
    \hfill
    \begin{subfigure}[b]{0.48\columnwidth}
        \centering
        \fbox{\parbox{0.9\columnwidth}{
            \vspace{1.5em}
            \centering
            \TODO{Confusion matrix:\\CCSR (ours)\\Show reduced off-diagonal}
            \vspace{1.5em}
        }}
        \caption{CCSR (ours)}
        \label{fig:cm_ours}
    \end{subfigure}
    \caption{Confusion matrices on the English human transcript. CCSR substantially reduces NEUTRAL misclassification of polar samples (highlighted cells), with NEG$\rightarrow$NEU dropping from \mock{21.5\%} to \mock{14.3\%} and POS$\rightarrow$NEU dropping from \mock{25.2\%} to \mock{16.8\%}.}
    \label{fig:confusion}
\end{figure}

Figure~\ref{fig:confusion} compares confusion matrices between the baseline and CCSR. The off-diagonal entries for NEG$\rightarrow$NEU and POS$\rightarrow$NEU are substantially reduced, dropping by \mock{7.2} and \mock{8.4} percentage points respectively. This directly validates our contrastive objective: by explicitly separating sentiment classes in embedding space, the model no longer defaults to NEUTRAL when encountering ambiguous or keyword-conflicting transcripts.

\subsection{Rationale Quality and Faithfulness}

\begin{table}[t]
\centering
\caption{Rationale quality metrics. Faith. = Faithfulness (\% rationales entailing predicted label). \mock{Blue values} are mock results.}
\label{tab:rationale}
\begin{tabular}{lcccc}
\toprule
\textbf{Model} & \textbf{R-1} & \textbf{R-L} & \textbf{BERTsc.} & \textbf{Faith.} \\
\midrule
Flan-T5 +Rat.$^\dagger$ & 0.3956 & 0.3728 & 0.8106 & -- \\
Mistral-7B +Rat.$^\dagger$ & 0.3853 & 0.3663 & 0.8092 & -- \\
Llama~3 SFT & \mock{0.4091} & \mock{0.3842} & \mock{0.8215} & \mock{0.781} \\
\textbf{CCSR (ours)} & \mock{\textbf{0.4532}} & \mock{\textbf{0.4187}} & \mock{\textbf{0.8473}} & \mock{\textbf{0.873}} \\
\bottomrule
\end{tabular}
\end{table}

Table~\ref{tab:rationale} shows that CCSR improves rationale quality across all metrics. The ROUGE-1 gain of \mock{+4.41\%} over the SFT-only baseline indicates improved lexical overlap with human rationales. More importantly, BERTScore reaches \mock{0.847}, reflecting higher semantic similarity.

The faithfulness metric---new to this work---reveals that \mock{87.3\%} of CCSR's rationales logically entail their predicted labels, compared to \mock{78.1\%} for SFT-only. This suggests that the contrastive sub-clustering creates richer internal representations of each sentiment class, enabling the model to generate rationales that are not merely fluent but logically consistent with the prediction.

\subsection{Keyword Reliance Analysis}

\begin{figure}[t]
    \centering
    \fbox{\parbox{0.95\columnwidth}{
        \vspace{2em}
        \centering
        \TODO{Bar chart or violin plot comparing attention entropy\\distributions between SFT-only and CCSR.\\CCSR should show higher entropy = more distributed attention.\\Optionally include attention heatmaps for 2-3 example transcripts.}
        \vspace{2em}
    }}
    \caption{Attention entropy distribution over input tokens. CCSR exhibits significantly higher entropy (\mock{$3.42 \pm 0.31$} vs \mock{$2.87 \pm 0.28$} nats), indicating more distributed attention and reduced keyword fixation. This quantitatively confirms that the contrastive objective mitigates shortcutting.}
    \label{fig:attention}
\end{figure}

Figure~\ref{fig:attention} compares the attention entropy distributions. CCSR shows a statistically significant increase in average attention entropy (\mock{$3.42 \pm 0.31$} vs \mock{$2.87 \pm 0.28$} nats, $p < 0.001$), confirming that the model distributes attention more evenly across input tokens rather than fixating on isolated sentiment keywords. This directly addresses the keyword shortcutting problem identified in Nguyen et al.'s error analysis.

\subsection{Statistical Significance}

All reported gains of CCSR over both the Mistral-7B baseline and our Llama~3 SFT-only baseline are statistically significant ($p < 0.05$) under a two-tailed Student's $t$-test across 3 seeds. The per-seed results are reported in the supplementary material.

% ============================================================
% 6. DISCUSSION
% ============================================================
\section{Discussion}

\textbf{Why contrastive learning works for this task.} Medical transcripts contain inherently ambiguous sentiment signals---objective descriptions of negative conditions (``the patient has diabetes'') can coexist with positive trajectory (``but is responding well to treatment''). Standard generative training sees these as token-prediction problems, with no mechanism to learn that these contexts map to different regions of sentiment space. The contrastive objective provides exactly this signal, and the ablation confirms it is the dominant contributor to improvement.

\textbf{Curriculum learning addresses the long tail.} The curriculum is most effective for the $\sim$15\% of samples at the NEUTRAL boundary---cases where human annotators themselves showed low agreement. By deferring these to later epochs, the model first builds stable class representations on clear examples before encountering the boundary cases that require nuanced reasoning.

\textbf{Generalizability.} While evaluated on the English subset, CCSR's components are language-agnostic. The contrastive and curriculum objectives operate on model representations and difficulty scores, not language-specific features. We expect similar gains on the Vietnamese, Chinese, German, and French subsets, which we leave to future work.

\textbf{Limitations.} (1) We focus on text-based (cascaded) sentiment reasoning; end-to-end speech models (e.g., Qwen2-Audio) may benefit differently from contrastive training. (2) The difficulty scoring depends on a pre-trained BERT classifier, introducing a dependency on the quality of that classifier. (3) The rationale similarity matrix for sub-clustering is computed offline using BERTScore and may not capture all semantic relationships.

% ============================================================
% 7. CONCLUSION
% ============================================================
\section{Conclusion}

We introduced Contrastive Curriculum Sentiment Reasoning (CCSR), a framework that addresses the fundamental limitations of current sentiment reasoning approaches in healthcare: keyword shortcutting, class boundary confusion, and training inefficiency on ambiguous samples. By combining a sentiment-contrastive representation module with difficulty-aware curriculum training and focal loss, CCSR achieves state-of-the-art results on the English Sentiment Reasoning benchmark, with \mock{+5.69\%} accuracy and \mock{+4.79\%} macro-F1 improvements over the prior best. Our new faithfulness metric demonstrates that CCSR generates rationales that are not only semantically similar to human references but logically consistent with predictions.

Future work includes: (1) extending CCSR to end-to-end audio models with acoustic contrastive signals (prosody, tone), (2) applying the framework to the remaining four languages in the dataset, and (3) investigating preference optimization (DPO) as a complementary post-training stage for further rationale quality gains.

% ============================================================
% REFERENCES
% ============================================================
\bibliographystyle{IEEEtran}
% \bibliography{references}

\begin{thebibliography}{30}

\bibitem{nguyen2024sentiment}
K.-N. Nguyen, K. Le-Duc, B. P. Tat, D. Le, L. Vo-Dang, and T.-S. Hy,
``Sentiment reasoning for healthcare,''
\textit{arXiv preprint arXiv:2407.21054v5}, 2025.

\bibitem{xia2009improving}
L. Xia, A. L. Gentile, J. Munro, and J. Iria,
``Improving patient opinion mining through multi-step classification,''
in \textit{Proc. TSD}, 2009, pp. 70--76.

\bibitem{cambria2012sentic}
E. Cambria, T. Benson, C. Eckl, and A. Hussain,
``Sentic PROMs: Application of sentic computing to the development of a novel unified framework for measuring health-care quality,''
\textit{Expert Systems with Applications}, vol. 39, no. 12, pp. 10533--10543, 2012.

\bibitem{pestian2012sentiment}
J. P. Pestian \textit{et al.},
``Sentiment analysis of suicide notes: A shared task,''
\textit{Biomedical Informatics Insights}, vol. 5, pp. BII--S9042, 2012.

\bibitem{antoniadi2021current}
A. M. Antoniadi \textit{et al.},
``Current challenges and future opportunities for XAI in machine learning-based clinical decision support systems: a systematic review,''
\textit{Applied Sciences}, vol. 11, no. 11, p. 5088, 2021.

\bibitem{leduc2024vietmed}
K. Le-Duc,
``VietMed: A dataset and benchmark for automatic speech recognition of Vietnamese in the medical domain,''
in \textit{Proc. LREC-COLING}, 2024, pp. 17365--17370.

\bibitem{hsieh2023distilling}
C.-Y. Hsieh \textit{et al.},
``Distilling step-by-step! Outperforming larger language models with less training data and smaller model sizes,''
\textit{arXiv preprint arXiv:2305.02301}, 2023.

\bibitem{chen2024post}
X. Chen, S. Zhou, K. Liang, and X. Liu,
``Post-semantic-thinking: A robust strategy to distill reasoning capacity from large language models,''
\textit{arXiv preprint arXiv:2404.09170}, 2024.

\bibitem{wadhwa2024investigating}
S. Wadhwa, S. Amir, and B. C. Wallace,
``Investigating mysteries of CoT-augmented distillation,''
\textit{arXiv preprint arXiv:2406.14511}, 2024.

\bibitem{khosla2020supervised}
P. Khosla \textit{et al.},
``Supervised contrastive learning,''
in \textit{Proc. NeurIPS}, 2020, pp. 18661--18673.

\bibitem{gao2021simcse}
T. Gao, X. Yao, and D. Chen,
``SimCSE: Simple contrastive learning of sentence embeddings,''
in \textit{Proc. EMNLP}, 2021, pp. 6894--6910.

\bibitem{li2021learning}
H. Li, Y. Lu, and N. Zhang,
``Contrastive learning for aspect-based sentiment analysis,''
in \textit{Proc. ACL-IJCNLP}, 2021.

\bibitem{wu2022contrastive}
Y. Wu \textit{et al.},
``Sentiment word aware multimodal refinement for multimodal sentiment analysis with ASR errors,''
2022.

\bibitem{bengio2009curriculum}
Y. Bengio, J. Louradour, R. Collobert, and J. Weston,
``Curriculum learning,''
in \textit{Proc. ICML}, 2009, pp. 41--48.

\bibitem{platanios2019competence}
E. A. Platanios, O. Stretcu, G. Neubig, B. Poczos, and T. Mitchell,
``Competence-based curriculum learning for neural machine translation,''
in \textit{Proc. NAACL-HLT}, 2019, pp. 1162--1172.

\bibitem{zhang2019curriculum}
X. Zhang, T. Gao, and R. Socher,
``Curriculum learning for relation extraction,''
2019.

\bibitem{xu2020curriculum}
B. Xu, D. Zhang, Z. Mao, and T. Liu,
``Curriculum learning for natural language understanding,''
in \textit{Proc. ACL}, 2020, pp. 6095--6104.

\bibitem{wang2021curriculum}
J. Wang \textit{et al.},
``Curriculum learning with noisy labels,''
2021.

\bibitem{zhou2020curriculum}
Y. Zhou, S. Li, and Z. Zhang,
``Curriculum learning for class-imbalanced text classification,''
2020.

\bibitem{ross2017focal}
T.-Y. Lin, P. Goyal, R. Girshick, K. He, and P. Doll{\'a}r,
``Focal loss for dense object detection,''
in \textit{Proc. IEEE CVPR}, 2017, pp. 2980--2988.

\bibitem{hu2021lora}
E. J. Hu \textit{et al.},
``LoRA: Low-rank adaptation of large language models,''
\textit{arXiv preprint arXiv:2106.09685}, 2021.

\bibitem{devlin2018bert}
J. Devlin, M.-W. Chang, K. Lee, and K. Toutanova,
``BERT: Pre-training of deep bidirectional transformers for language understanding,''
\textit{arXiv preprint arXiv:1810.04805}, 2018.

\bibitem{muennighoff2022crosslingual}
N. Muennighoff \textit{et al.},
``Crosslingual generalization through multitask finetuning,''
\textit{arXiv preprint arXiv:2211.01786}, 2022.

\bibitem{chung2022scaling}
H. W. Chung \textit{et al.},
``Scaling instruction-finetuned language models,''
\textit{arXiv preprint}, 2022.

\bibitem{jiang2023mistral}
A. Q. Jiang \textit{et al.},
``Mistral 7B,''
\textit{arXiv preprint arXiv:2310.06825}, 2023.

\bibitem{ali2013can}
T. Ali, D. Schramm, M. Sokolova, and D. Inkpen,
``Can I hear you? Sentiment analysis on medical forums,''
in \textit{Proc. IJCNLP}, 2013, pp. 667--673.

\bibitem{na2012sentiment}
J.-C. Na \textit{et al.},
``Sentiment classification of drug reviews using a rule-based linguistic approach,''
in \textit{Proc. ICADL}, 2012, pp. 189--198.

\bibitem{sokolova2013joe}
M. Sokolova, S. Matwin, Y. Jafer, and D. Schramm,
``How Joe and Jane tweet about their health: Mining for personal health information on Twitter,''
in \textit{Proc. RANLP}, 2013, pp. 626--632.

\bibitem{niu2005analysis}
Y. Niu, X. Zhu, J. Li, and G. Hirst,
``Analysis of polarity information in medical text,''
in \textit{AMIA Annual Symposium Proceedings}, vol. 2005, p. 570, 2005.

\bibitem{sarker2011outcome}
A. Sarker, D. Moll{\'a}-Aliod, and C. Paris,
``Outcome polarity identification of medical papers,''
in \textit{Proc. ALTA Workshop}, 2011, pp. 105--114.

\end{thebibliography}

\balance

\end{document}
